\section{Related Work}

ZeroFalse and QASecClaw are the two systems most closely related to Polygraph.
Both combine a \ac{sast} scanner with an \ac{llm} that reviews each finding to
decide whether it is a real vulnerability, the same basis Polygraph builds on,
and both decide it in a single pass. Xiong and Zhang instead evaluate autonomous
agents that investigate every finding on their own, which is the second stage
Polygraph applies selectively when necessary.

ZeroFalse \cite{iranmanesh2025zerofalseimprovingprecisionstatic} operates only
on findings from CodeQL. It treats such alerts as structured contracts and
enriches them with the scanner's own data-flow trace, the surrounding code, and
information specific to each \ac{cwe} class before a
single \ac{llm} pass returns a verdict. Across the OWASP Java Benchmark and a curated
real-world dataset it reaches F1-scores of 0.912 and 0.955, below Polygraph's
0.977 on the Python variant of the OWASP Benchmark. Its design is based
on \ac{cwe}-specific tailoring, and the authors name contextual blindness, the
difficulty of supplying sufficient context for a multi-file finding, as an open
problem. Polygraph addresses this problem with the autonomous escalation agent that gathers evidence on-demand.
It also maintains no prompt per weakness class, but uses a generic
prompt template for all findings that includes a description of the concrete
finding. Furthermore, rather than reuse the data-flow traces the scanner potentially
emitted, Polygraph computes its own taint analysis on-demand,
which removes the dependence on a scanner that populates \ac{sarif} flow traces
comprehensively. ZeroFalse depends on CodeQL for exactly this reason. As a
result, it cannot triage findings of arbitrary scanners, while Polygraph sits
between any \ac{sarif}-compatible scanner and viewer.

QASecClaw \cite{ameen2026qasecclawmultiagentllmapproach} orchestrates several
role-specific agents that cover scan planning, Semgrep-based scanning, evidence
correlation, and an \ac{llm} \ac{sast} Filter Agent. It classifies each finding from its
enclosing source file and \ac{cwe} label, with a fail-open policy that retains a
finding whenever the model fails. On OWASP Benchmark v1.2 it achieves an F1-score
of 0.909 and removes 88.6\% of \acp{fp}. Its filter, however,
performs no data-flow analysis and reasons only over the enclosing file. Its
agents also form a fixed pipeline rather than an investigator that gathers
missing evidence when needed. In contrast, on the Python variant of the OWASP Benchmark,
Polygraph reaches an F1-score of 0.977 while suppressing 95.8\% of the \ac{fp} traps.

Xiong and Zhang \cite{xiong2026} evaluate three general-purpose coding agents,
Aider, OpenHands, and SWE-agent, as \ac{fp} filters under three different
models, rather than proposing a system of their own. Each agent receives the
alert and the source code and settles the finding on its own. On OWASP Benchmark
v1.2 for Java they reduce an initial \ac{fp} rate of 98.3\% to 6.3\%, and the
result carries over to C/C++ findings.

ZeroFalse and QASecClaw judge every finding in a single, fixed-context pass, and
neither recovers when the evidence a verdict needs is missing from that context.
Xiong and Zhang measure what this costs. A baseline that receives exactly the
files their agent had read, but answers in one pass, reaches an F1-score of
0.516 against the agent's 0.955, so the deciding factor is the iterative
gathering of evidence rather than a larger context. Polygraph therefore treats
uncertain as a first-class verdict. The cheap first stage answers the findings
it can, and only those it cannot resolve escalate to an autonomous agent that
explores the source code. This confines expensive exploration to the
hard-to-decide cases that the two prior systems leave unresolved, directly
attacks the contextual-blindness limitation that ZeroFalse names, and costs far
less than running an agent on every finding, at 8,196 tokens per OWASP finding on average
against the 130,846 tokens per alert their best configuration needs.

The same design also decides how many real vulnerabilities survive. Their best
configuration classifies 22.25\% of the benchmark's real vulnerabilities as
\acp{fp}, which leads the authors to reject such agents for automatic suppression. Polygraph reaches a
dangerous-suppression rate of 0.4\% on the Python variant of the OWASP benchmark, a different corpus
but the same weakness families, because its escalation trigger is asymmetric. It
fires on every uncertain verdict and every \ac{fp} below high certainty, but
never on a \ac{tp}, so suppression is the only verdict that has to survive a
second pass.

Two further differences widen Polygraph's applicability. All three evaluations
cover Java, ZeroFalse through CodeQL, QASecClaw through Semgrep, and Xiong and
Zhang through four scanners at once and a small C/C++ sample. In contrast,
Polygraph supports Python, C, C++, JavaScript, and TypeScript first-class, but can
also handle findings in other languages. Where ZeroFalse consumes
\ac{sarif} only from CodeQL and QASecClaw consumes Semgrep JSON, Polygraph both reads and
writes regular \ac{sarif}, marking \acp{fp} through the standard
\texttt{suppressions} field. This way, it slots between any compliant scanner
and viewer. A prompt- and tool-hashing cache further makes repeated scans of an
unchanged codebase effectively free, a concern that none of the three addresses.

On their respective OWASP evaluations, ZeroFalse and QASecClaw report F1-scores
of 0.912 and 0.909, while Polygraph reaches 0.977. This gap has two causes.
Polygraph suppresses more \acp{fp} and, at the same time, misclassifies fewer
real vulnerabilities as \acp{fp}, so less noise reaches the reviewer without
more real vulnerabilities being hidden. The agents of Xiong and Zhang achieve
only the first, as they remove a comparable share of the noise but hide far more
real vulnerabilities. Only suppression that is safe in both directions can be
trusted, and that is what lets a team act on Polygraph's verdicts and re-examine
the escalated cases instead of double-checking the whole report.
